\documentclass[9pt]{article}
% Kiln CV template: a white drawing sheet, IBM Plex, ink rules, one ember thing (the mark).
% The alternative Kiln-style CV (kept for later; the main CV is cv/CV_Borodin_EN.tex). Shared by cv/kiln/CV_Borodin_EN.tex and CV_Borodin_RU.tex. Build from the repo root with
%   xelatex -output-directory=cv/build cv/kiln/CV_Borodin_EN.tex
% Fonts: IBM Plex Sans Condensed and IBM Plex Mono from IBM's own release (Cyrillic included).

\usepackage[a4paper,top=10mm,bottom=7mm,left=12mm,right=12mm]{geometry}
\usepackage{fontspec}
\usepackage{xcolor}
\usepackage{graphicx}
\usepackage{tikz}
\usepackage{enumitem}
\usepackage{tabularx}
\usepackage{array}
\usepackage[hidelinks,unicode]{hyperref}

% ---------------------------------------------------------------- tokens (Kiln 4 · Colour)
\definecolor{ink}{HTML}{16130F}
\definecolor{ink2}{HTML}{605A52}
\definecolor{dimmed}{HTML}{9F998F}
\definecolor{rule}{HTML}{D1CDC6}
\definecolor{ember}{HTML}{E84C14}

% ---------------------------------------------------------------- type (Kiln 3 · Type)
\setmainfont{IBM Plex Sans Cond}[BoldFont={IBM Plex Sans Cond SmBld}, ItalicFont={IBM Plex Sans Cond}]
\newfontfamily\monofont{IBM Plex Mono Medm}[BoldFont={IBM Plex Mono SmBld}]
\newfontfamily\labelfont{IBM Plex Mono Medm}[LetterSpace=11]

\color{ink}
\pagestyle{empty}
\setlength{\parindent}{0pt}
\setlength{\parskip}{0pt}
\setlength{\tabcolsep}{0pt}
\emergencystretch=2em
\raggedright

\hyphenpenalty=10000 \exhyphenpenalty=10000   % ragged right, no broken words (Kiln 3 · Type)

% the label column, used where a label names a row (skills)
\newlength{\labw}\setlength{\labw}{25mm}
\newlength{\bodyw}\setlength{\bodyw}{\dimexpr\textwidth-\labw\relax}

% ---------------------------------------------------------------- pieces
% Mono capitals with +0.11em, the label voice
\newcommand{\lab}[1]{{\labelfont\fontsize{6.6}{9}\selectfont\color{ink2}\MakeUppercase{#1}}}
% Dates and small technical facts in mono
\newcommand{\mono}[1]{{\monofont\fontsize{7.4}{10}\selectfont\color{ink2}#1}}
% A link shown as its address, quiet
\newcommand{\addr}[2]{\href{#1}{\mono{#2}}}

% Section: a heavy ink rule, then the label (the site's section pattern)
\newcommand{\cvsection}[1]{%
  \par\vspace{3pt}%
  {\color{ink}\rule{\textwidth}{1.4pt}}\par
  \vspace{2.5pt}\lab{#1}\par\vspace{3pt}%
}

% An entry: title left, date right in mono (the site's job row), place under it, then the list
\newcommand{\entry}[4]{% #1 date, #2 title, #3 place, #4 list
  \par\noindent{\bfseries\fontsize{9.8}{12}\selectfont #2}\hfill\mono{#1}\par
  {\color{ink2}#3}\par
  #4\par\vspace{3pt}%
}

% A labelled row: mono label left, text right (\leavevmode keeps both first baselines aligned)
\newcommand{\row}[2]{%
  \par\noindent
  \parbox[t]{\labw}{\leavevmode\raggedright\mono{#1}}%
  \parbox[t]{\dimexpr\linewidth-\labw\relax}{\leavevmode\raggedright #2}\par\vspace{2.5pt}%
}

% A one-line fact with its date on the right
\newcommand{\fact}[2]{\par\noindent #1\hfill\mono{#2}\par\vspace{1pt}}

\setlist[itemize]{leftmargin=1em,labelsep=0.35em,itemsep=0pt,topsep=1.5pt,parsep=0pt,partopsep=0pt}
\renewcommand\labelitemi{{\color{dimmed}—}}

% The photo, cropped square with Kiln's panel radius
\newcommand{\photo}[2]{%
  \begin{tikzpicture}
    \clip[rounded corners=2.2mm] (0,0) rectangle (#2,#2);
    \node[anchor=south west,inner sep=0] at (0,0) {\includegraphics[width=#2,height=#2]{#1}};
  \end{tikzpicture}%
}

% Header: lockup left, contacts right, the heavy rule under it (the site header)
\newcommand{\cvheader}[2]{% #1 lockup pdf, #2 contact lines
  \noindent
  \begin{minipage}[c]{0.5\textwidth}\includegraphics[height=11mm]{#1}\end{minipage}%
  \begin{minipage}[c]{0.5\textwidth}\raggedleft #2\end{minipage}\par
  \vspace{5pt}{\color{ink}\rule{\textwidth}{1.4pt}}\par\vspace{7pt}%
}

% Lede size (Russian sets wider, so its file may override this)
\newcommand{\ledesize}{\fontsize{10}{12.6}\selectfont}

% Title block: label, name, lede; photo on the right
\newcommand{\cvtitle}[4]{% #1 label, #2 name, #3 lede, #4 photo file
  \noindent
  \begin{minipage}[t]{\dimexpr\textwidth-27mm\relax}
    \raggedright\lab{#1}\par\vspace{3pt}
    {\bfseries\fontsize{22}{25}\selectfont #2}\par\vspace{3pt}
    {\ledesize\color{ink2}#3}
  \end{minipage}\hfill
  \begin{minipage}[t]{22mm}\vspace{0pt}\photo{#4}{22mm}\end{minipage}\par
}

\fontsize{9}{11}\selectfont

% Two columns of labelled rows (short values, e.g. skills)
\newenvironment{twocol}{\par\noindent\begin{minipage}[t]{0.485\textwidth}}{\end{minipage}\par}
\newcommand{\nextcol}{\end{minipage}\hfill\begin{minipage}[t]{0.485\textwidth}}

\hypersetup{
  pdftitle={Александр Бородин -- ML-инженер | LLM и AI-агенты},
  pdfauthor={Александр Бородин},
  pdfsubject={Резюме}
}

% Cyrillic sets wider: a quarter point smaller keeps the CV on one page
\renewcommand{\ledesize}{\fontsize{9.6}{12.2}\selectfont}
\setlength{\labw}{21mm}   % short Russian labels: give the values the width

\begin{document}
\fontsize{8.75}{10.5}\selectfont

\cvheader{cv/kiln/lockup-ru.pdf}{%
  \href{mailto:apdorodin765@gmail.com}{\mono{apdorodin765@gmail.com}}\quad
  \mono{+7 908 971 16 97}\par
  \addr{https://github.com/aleksanderborodin}{github.com/aleksanderborodin}\quad
  \addr{https://t.me/aleksanderbor}{t.me/aleksanderbor}\par
  \addr{https://aleksanderborodin.github.io/ru/}{aleksanderborodin.github.io/ru}}

\cvtitle{ML-инженер · LLM и AI-агенты}{Александр Бородин}%
{ML-инженер, студент ФПМИ МФТИ. Разрабатываю системы на~основе LLM: от~RAG и~интеграции моделей до~мультиагентных систем и~агентной эволюции~— итеративного поиска и~улучшения решений с~помощью агентов. Есть опыт построения бенчмарков, оценки моделей и~решения задач оптимизации.}%
{assets/profile-original.jpg}

\cvsection{Опыт работы}
\entry{фев–авг 2026 · 6 мес}{ML-инженер~— исследовательская практика}{ЦКМ МФТИ}{%
\begin{itemize}
  \item Разрабатывал агентов на~основе LLM и~бенчмарки для оценки качества их работы.
  \item Работал над агентной эволюцией: системами, в~которых агенты итеративно генерируют, проверяют и~улучшают решения.
  \item Участвовал в~подготовке двух научных статей: одна находится на~рецензировании в~ACL, вторая готовится к~подаче.
\end{itemize}}
\entry{янв–май 2026}{Product / AI Engineer}{\href{https://modelgate.ru}{ModelGate.ru}~— LLM API-шлюз и~прокси-сервис}{%
\begin{itemize}
  \item Разрабатывал LLM API-шлюз с~единым интерфейсом для доступа к~моделям разных провайдеров.
  \item Интегрировал API провайдеров и~реализовывал маршрутизацию запросов с~выбором модели и~провайдера.
  \item Создавал основные компоненты сервиса, включая backend и~публичный сайт. Стек: Python, LLM API.
\end{itemize}}
\entry{сен 2024 – фев 2025}{Разработчик моделей и~алгоритмов}{Розничный банк, клиентский проект}{%
\begin{itemize}
  \item Разработал алгоритм оптимизации инкассации; анализ исторических данных дал прогноз +20~млн~руб. годовой прибыли.
  \item Построил пайплайн для скрапинга и~анализа данных ЦБ РФ, недоступных через API, и~внутреннюю RAG-систему для вопросов по~документам.
  \item Стек: Python, NumPy, Pandas, RAG, скрапинг.
\end{itemize}}

\cvsection{Проекты}
\entry{2026}{Idea Evolve}{\addr{https://github.com/aleksanderborodin/idea_evolve}{github.com/aleksanderborodin/idea\_evolve}}{%
\begin{itemize}
  \item Эволюционный мультиагентный фреймворк: специализированные Claude-агенты работают параллельно по~поколениям и~строят общую базу знаний.
  \item Улучшил поиск множества Сидона с~66 до~105 элементов за~7 автономных поколений, выше ориентира √N~≈~100 при N~=~10\,000.
  \item Реализовал ролевую специализацию агентов, возобновляемое состояние, извлечение знаний и~Flask-дашборд.
  \item Стек: Python 3.12+, Claude Code CLI, OpenCode, Anthropic API, Flask, мультиагентная оркестрация.
\end{itemize}}
\entry{2025}{Prefix Injection Jailbreak}{\addr{https://github.com/aleksanderborodin/prefix_injection_jailbreak}{github.com/aleksanderborodin/prefix\_injection\_jailbreak}}{%
\begin{itemize}
  \item Исследовал префиксные атаки, вдохновлённые Choice Blindness: шаблоны заставляли Gemini 2.5 Pro и~DeepSeek R1 завершать MaliciousInstruct-запросы почти в~100\% случаев против почти нулевого базового уровня.
  \item Стек: Python, LLM API, оценка моделей, проектирование промптов и~контекста.
\end{itemize}}
\entry{2025}{Yoola Explain}{\addr{https://github.com/aleksanderborodin/yoola_explain}{github.com/aleksanderborodin/yoola\_explain}}{%
\begin{itemize}
  \item Браузерное расширение: извлекает Terms of Service и~возвращает структурированное LLM-резюме~— ключевые пункты, сбор данных и~предупреждения.
  \item Добавил мультиязычность и~кэширование. Стек: JavaScript, FastAPI, OpenRouter, LLM API.
\end{itemize}}

\cvsection{Образование}
\entry{2023–2027}{Бакалавриат, прикладная математика и~физика}{ФПМИ МФТИ, учусь · GPA 7.8/10}{%
Курсы: матанализ, теория вероятностей, линейная алгебра, алгоритмы, вычислительная сложность, ML, физика, экономика.}

\cvsection{Навыки}
\begin{twocol}
\row{Код}{Python: продвинутый; C++: средний; чтение кода на~C и~Assembler.}
\row{Данные / ML}{PyTorch: базовый; NumPy; Pandas; PostgreSQL; Matplotlib; ClearML; Excel.}
\row{NLP / LLM}{BM25; эмбеддинги; RAG; LLM API; оценка моделей; проектирование промптов и~контекста; тестирование джейлбрейков.}
\nextcol
\row{Агенты}{Claude Code CLI; OpenCode; мультиагентная оркестрация; файловые процессы.}
\row{Бэкенд}{FastAPI; Flask; LangChain; LangGraph; Linux: домашний сервер и~рабочая станция.}
\row{Языки}{Русский: родной; английский: B2+ по~placement test МФТИ.}
\end{twocol}

\cvsection{Образовательные программы и достижения}
\entry{21 июл – 4 авг 2026}{\href{https://airi.net/ru/summer-school-2026/}{«Лето с~AIRI 2026»}~— летняя школа по~искусственному интеллекту}{Итоговый командный проект: \href{https://github.com/n4rly-boop/ideas-lake}{«Озеро идей»}~— долговременная память для эволюционных AI-систем с~гибридным поиском и~графом знаний.}{}
\fact{\textbf{Золотая медаль}, \href{https://iepho.ru/}{Международная экспериментальная олимпиада по~физике (IEPhO)}}{2023}
\fact{\textbf{Призёр заключительного этапа} \href{https://vos.olimpiada.ru/}{ВсОШ по~физике}~— 10 и~11 классы}{2022, 2023}

\end{document}
